\documentclass[10pt,a4paper]{amsart}
\usepackage[margin=19mm]{geometry}
\usepackage{amsmath,amssymb,mathtools}
\usepackage[scr=rsfs,cal=euler]{mathalfa}
\usepackage{xfrac}
\usepackage[unicode,hidelinks,bookmarks=false]{hyperref}
\newcommand{\ie}{i.e.\ }
\newcommand{\cf}{cf.\ }
\newcommand{\resp}{respectively\ }
\newcommand{\Subsection}{\subsection}
\providecommand{\nobreakdash}{\nobreak\hbox{-}}
\setlength{\emergencystretch}{2em}
\multlinegap0pt
\usepackage{amssymb}
\usepackage{enumerate}
\newtheorem{theorem}{Theorem}
\title{A representation for the integral kernel of the composition of multivariate Bernstein-Durrmeyer operators}
\author{Ulrich Abel, Ana Maria Acu, Margareta Heilmann, Ioan Rasa}
\date{Source: arXiv:2512.17297v1 (CC BY 4.0)}
\begin{document}
\maketitle

\noindent\emph{Source and licence.} A representation for the integral kernel of the composition of multivariate Bernstein-Durrmeyer operators. Version arXiv:2512.17297v1 (CC BY 4.0), distributed by arXiv under the Creative Commons Attribution 4.0 International licence, http://creativecommons.org/licenses/by/4.0/, as recorded in the arXiv licence metadata for this exact version and shown on the abstract page https://arxiv.org/abs/2512.17297v1. Full licence notice: \texttt{licenses/arxiv-2512.17297-CC-BY-4.0.txt}.\par\smallskip

\noindent\textit{Editorial scope.} Counted unit: the article's theorem theorem (source0.tex lines 289--298). The article's complete original proof of it is delivered with it (source0.tex lines 300--397, one \texttt{proof} environment of 98 lines). No mathematics is written, repaired or re-derived: the statement and the proof are the article's own lines, in the article's own notation, and no retained line is substituted or re-flowed.  No further source-local context is needed: every cross-reference of every retained line resolves inside the two delivered spans.  Nothing was omitted from any retained span, so the package carries no omission record.  The retained lines cite 1 works, whose entries are transcribed verbatim from the article's own bibliography source in the e-print and delivered with short editorial labels recorded in the item's reference mapping.  No retained line contains a figure environment, an image-inclusion command or drawing code, so no image is delivered and the package declares no asset dependency.  Editorial formatting only, in addition: an AMS \texttt{amsart} preamble that restates the article's own declarations and macros used by the retained lines, namely the environments \texttt{theorem} on separate counters, together with the standard packages those lines need.  The retained environments print in this document as Theorem 1 in order of appearance, whereas the article prints the same items with the numbering of its own full document.  The complete native source of the article as distributed in the arXiv e-print for this exact version is bundled unchanged as \texttt{sources/191307/source0.tex}.
\begin{theorem}
For positive integers $m,n$ and $x,y\in \mathbb{R}^{d}$, the kernel $%
K_{m,n}\left( x,y\right) $ has the representation 
\begin{equation*}
K_{m,n}\left( x,y\right) =\frac{\left( m+d\right) !\left( n+d\right) !}{%
\left( m+n+d\right) !}\sum_{\ell \geq 0}\binom{m}{\left \vert \ell \right
\vert }\binom{n}{\left \vert \ell \right \vert }B_{\ell }\left( \mathbf{x}%
\right) B_{\ell }\left( \mathbf{y}\right) .
\end{equation*}
\end{theorem}

\begin{proof}
For positive integers $m,n$ and $x,y\in \mathbb{R}^{d}$, we have 
\begin{equation*}
K_{m,n}\left( x,y\right) =\sum_{\left\vert \beta \right\vert
=m}\sum_{\left\vert \alpha \right\vert =n}B_{\alpha }\left( y\right)
B_{\beta }\left( x\right) \frac{1}{\left\langle \mathbf{1},B_{\alpha
}\right\rangle \left\langle \mathbf{1},B_{\beta }\right\rangle }\int_{%
\mathbb{S}^{d}}B_{\alpha }\left( t\right) B_{\beta }\left( t\right) dt.
\end{equation*}%
Note that $\left\langle \mathbf{1},B_{\alpha }\right\rangle =\left\vert
\alpha \right\vert !/\left( \left\vert \alpha \right\vert +d\right) !$ (see,
e.g., \cite[Page~147]{Berens-Schmidt-Xu-JAT-1992}). Since 
\begin{equation*}
\int_{S^{d}}B_{\alpha }\left( t\right) B_{\beta }\left( t\right) dt=\frac{%
\binom{\left\vert \alpha \right\vert }{\alpha }\binom{\left\vert \beta
\right\vert }{\beta }}{\binom{\left\vert \alpha +\beta \right\vert }{\alpha
+\beta }}\left\langle \mathbf{1},B_{\alpha +\beta }\right\rangle 
\end{equation*}%
we obtain 
\begin{eqnarray*}
K_{m,n}\left( x,y\right)  &=&\frac{\left( m+d\right) !\left( n+d\right)
!\left( m+n\right) !}{m!n!\left( m+n+d\right) !}\sum_{\left\vert \beta
\right\vert =m}\sum_{\left\vert \alpha \right\vert =n}B_{\alpha }\left(
y\right) B_{\beta }\left( x\right) \frac{\binom{\left\vert \alpha
\right\vert }{\alpha }\binom{\left\vert \beta \right\vert }{\beta }}{\binom{%
\left\vert \alpha +\beta \right\vert }{\alpha +\beta }} \\
&=&\frac{\left( m+d\right) !\left( n+d\right) !}{\left( m+n+d\right) !}%
\sum_{\left\vert \beta \right\vert =m}\frac{1}{\beta !}B_{\beta }\left(
x\right) \sum_{\left\vert \alpha \right\vert =n}B_{\alpha }\left( y\right) 
\frac{\left( \alpha +\beta \right) !}{\alpha !}.
\end{eqnarray*}%
With the notation\ $\mathbf{t}=\left( t_{0},t_{1},\ldots ,t_{d}\right) $\
and $\mathbf{D}^{\beta }=\prod_{\nu =0}^{d}\left( \frac{\partial }{\partial
t_{\nu }}\right) ^{\beta _{\nu }}$, the inner sum is equal to 
\begin{eqnarray*}
\sum_{\left\vert \alpha \right\vert =n}B_{\alpha }\left( y\right) \frac{%
\left( \alpha +\beta \right) !}{\alpha !} &=&\sum_{\left\vert \alpha
\right\vert =n}B_{\alpha }\left( y\right) \left( \mathbf{D}^{\beta }\mathbf{t%
}^{\alpha +\beta }\right) _{\mid \mathbf{t=1}}=\mathbf{D}^{\beta }\left( 
\mathbf{t}^{\beta }\sum_{\left\vert \alpha \right\vert =n}B_{\alpha }\left(
y\right) \mathbf{t}^{\alpha }\right) _{\mid \mathbf{t=1}} \\
&=&\left( \prod_{\nu =0}^{d}\left( \frac{\partial }{\partial t_{\nu }}%
\right) ^{\beta _{\nu }}\left( t_{0}^{\beta _{0}}t_{1}^{\beta _{1}}\cdots
t_{d}^{\beta _{d}}\left( y_{0}t_{0}+y_{1}t_{1}+\cdots +y_{d}t_{d}\right)
^{n}\right) \right) _{\mid \mathbf{t=1}} \\
&=&\sum_{\ell \leq \beta }n^{\underline{\left\vert \ell \right\vert }%
}\prod_{\nu =0}^{d}\left( \binom{\beta _{\nu }}{\ell _{\nu }}\frac{\beta
_{\nu }!}{\ell _{\nu }!}y_{\nu }^{\ell _{\nu }}\right)  \\
&=&\sum_{\ell \leq \beta }n^{\underline{\left\vert \ell \right\vert }}\frac{1%
}{\left\vert \ell \right\vert !}B_{\ell }\left( y\right) \beta !\prod_{\nu
=0}^{d}\binom{\beta _{\nu }}{\ell _{\nu }}.
\end{eqnarray*}%
Here, $\ell \leq \beta $ means summation over all nonnegative integers $\ell
_{\nu }$ such that $\ell _{\nu }\leq \beta _{\nu }$, for $\nu =0,\ldots ,d$.
Therefore, 
\begin{eqnarray*}
&&\left( \frac{\left( m+d\right) !\left( n+d\right) !}{\left( m+n+d\right) !}%
\right) ^{-1}K_{m,n}\left( x,y\right)  \\
&=&\sum_{\left\vert \beta \right\vert =m}B_{\beta }\left( x\right)
\sum_{\ell \leq \beta }\binom{n}{\left\vert \ell \right\vert }B_{\ell
}\left( y\right) \prod_{\nu =0}^{d}\binom{\beta _{\nu }}{\ell _{\nu }} \\
&=&\sum_{\ell \geq 0}\binom{n}{\left\vert \ell \right\vert }B_{\ell }\left(
y\right) \sum_{\left\vert \beta \right\vert =m}B_{\beta }\left( x\right)
\prod_{\nu =0}^{d}\binom{\beta _{\nu }}{\ell _{\nu }} \\
&=&\sum_{\ell \geq 0}\binom{n}{\left\vert \ell \right\vert }B_{\ell }\left(
y\right) \sum_{\left\vert \beta \right\vert =m}{\binom{\left\vert \beta
\right\vert }{\beta }}\mathbf{x}^{\beta }\prod_{\nu =0}^{d}\frac{\left(
\beta _{\nu }\right) ^{\underline{\ell _{\nu }}}}{\ell _{\nu }!} \\
&=&\sum_{\ell \geq 0}\binom{n}{\left\vert \ell \right\vert }B_{\ell }\left(
y\right) m^{\underline{\left\vert \ell \right\vert }}\left( \prod_{\nu
=0}^{d}\frac{1}{\ell _{\nu }!}\right) \sum_{\substack{ \left\vert \beta
\right\vert =m \\ \beta \geq \ell }}{\binom{m-\left\vert \ell \right\vert }{%
\beta -\ell }}\mathbf{x}^{\beta },
\end{eqnarray*}%
where $\ell \geq 0$ means summation over all nonnegative integers $\ell
_{\nu }$ $\left( \nu =0,\ldots ,d\right) $. Observing that 
\begin{equation*}
\sum_{\substack{ \left\vert \beta \right\vert =m \\ \beta \geq \ell }}{%
\binom{m-\left\vert \ell \right\vert }{\beta -\ell }}\mathbf{x}^{\beta
}=\sum_{\left\vert \beta \right\vert =m-\left\vert \ell \right\vert }{\binom{%
m-\left\vert \ell \right\vert }{\beta }}\mathbf{x}^{\beta +\ell }=\mathbf{x}%
^{\ell }\sum_{\left\vert \beta \right\vert =m-\left\vert \ell \right\vert }{%
\binom{m-\left\vert \ell \right\vert }{\beta }}B_{\beta }\left( x\right) =%
\mathbf{x}^{\ell }
\end{equation*}%
and using that $\prod_{\nu =0}^{d}\frac{1}{\ell _{\nu }!}=\frac{1}{%
\left\vert \ell \right\vert !}\binom{\left\vert \ell \right\vert }{\ell }$
we obtain 
\begin{equation*}
\left( \frac{\left( m+d\right) !\left( n+d\right) !}{\left( m+n+d\right) !}%
\right) ^{-1}K_{m,n}\left( x,y\right) =\sum_{\ell \geq 0}\binom{m}{%
\left\vert \ell \right\vert }\binom{n}{\left\vert \ell \right\vert }B_{\ell
}\left( y\right) \binom{\left\vert \ell \right\vert }{\ell }\mathbf{x}^{\ell
}.
\end{equation*}%
Since $\binom{\left\vert \ell \right\vert }{\ell }\mathbf{x}^{\ell }=B_{\ell
}\left( x\right) $ the proof is completed.
\end{proof}

\begin{thebibliography}{99}
\bibitem[Berens]{Berens-Schmidt-Xu-JAT-1992} H. Berens, H.J. Schmidt, and Y. Xu, Bernstein--Durrmeyer polynomials on a simplex, J. Approx. Theory \textbf{68}% , no. 3, 247--261, 1992. \newline https://doi.org/10.1016/0021-9045(92)90104-V
\end{thebibliography}
\end{document}
